\section{L2 -- Data Decomposition}

\subsection{Data Decomposition}

\mult{2}

\begin{definition}{Data Decomposition}\\
    Split the \emph{data structure} (image array) into chunks updated independently. Scales with \#chunks.
\end{definition}

\begin{concept}{Granularity}\\
    \begin{itemize}
        \item big chunk: low parallelism, low overhead
        \item small chunk: high parallelism, high overhead + all-boundary
        \item largest = whole image; smallest = 1 pixel
    \end{itemize}
\end{concept}

\multend

\begin{KR}{Recipe: Choose \& Size Data Chunks}\\
    \paragraph{Feasible?} heavy work = data-structure manipulation + similar op applied independently?
    \paragraph{Pattern} array + local neighbourhood $\to$ Geometric Decomposition.
    \paragraph{Shape} minimise surface/volume $\to$ square tiles.
    \paragraph{Boundaries} 1-cell ghost halo; exchange a-priori or overlapped.
    \paragraph{Map} \#chunks $>$ \#UEs for balancing.
\end{KR}

\begin{example2}{Ex.\ S6--8 -- Feasible for edge detection?}\\
    Is data decomposition feasible? Tasks? Granularity boundaries?
    \tcblower
    \important{Yes}: Sobel over every pixel, kernel needs only a $3\times3$ neighbourhood $\to$ regions independent. Tasks = per-region filters. Whole image = no parallelism; 1 pixel = overhead dominates.
\end{example2}

\subsection{Geometric Decomposition}

\begin{concept}{Geometric Decomposition}\\
    Array problems with local-neighbourhood (stencil) operations $\to$ split into tiles. Interior = embarrassingly parallel; only boundaries shared.
\end{concept}

\begin{definition}{Ghost / Halo Cells}\\
    Read-only copy of a neighbour's boundary, needed because the $3\times3$ kernel reads one row/col out. Exchange (1) a-priori or (2) overlapped with compute.
\end{definition}

\begin{formula}{Surface vs Volume}\\
    compute $\propto$ \emph{volume}; communication $\propto$ \emph{surface}.
    \vspace{1mm}\\
    equal volume $N^2/4$: square surface $2N$ vs strip $\tfrac{5N}{2}$ $\Rightarrow$ \important{square wins}.
    \vspace{1mm}\\
    $10\times10$: square 25 calc/20 transfers; strip 20/25.
    % INSERT IMAGE: L2 slide 13 -- surface vs volume sketch
\end{formula}

\begin{KR}{Recipe: Surface/Volume Comparison}\\
    Volume = cells (compute); Surface = boundary cells (comms). Lower surface/volume wins for equal volume.
\end{KR}

\begin{example2}{Ex.\ S12/13 -- Square vs rectangle}\\
    For a $10\times10$ matrix in 4 chunks, which shape is better and why?
    \tcblower
    Equal volume, but square surface $2N <$ strip $\tfrac{5N}{2}$ $\Rightarrow$ \important{square} = more compute per unit communication.
\end{example2}

\begin{example2}{Ex.\ S16/17 -- Apply to edge detection}\\
    Task parallelism or geometric? Chunks? Data sharing?
    \tcblower
    \important{Geometric Decomposition}; square tiles (or strips); share a 1-pixel \important{read-only ghost halo}; interior embarrassingly parallel.
\end{example2}

\subsection{Dependency \& Mapping}

\begin{concept}{Distributed Array}\\
    More chunks than UEs, assigned block-cyclic / round-robin so each UE owns several scattered chunks $\to$ smoother load balance.
\end{concept}

\begin{example2}{Ex.\ S10 -- Group \& order chunks}\\
    How to group and order chunk tasks?
    \tcblower
    Group tasks on the same chunk; order = front pipeline, then \important{ghost-exchange before boundary}, then \texttt{sobel\_x/y}$\to$\texttt{sobel\_all} join per chunk.
\end{example2}

\subsection{SPMD}

\begin{concept}{SPMD}\\
    \emph{Single Program Multiple Data}: one program on every UE, each with a unique ID/rank used to branch and index its data.
\end{concept}

\begin{KR}{Recipe: SPMD Program}\\
    Init $\to$ get ID/rank $\to$ common code (branch/index by ID) $\to$ distribute data $\to$ collect+combine $\to$ finalise.
\end{KR}

\begin{example2}{Ex.\ S27/28 -- SPMD evaluated}\\
    Strengths and limits?
    \tcblower
    One maintainable, portable code base; overhead at start-up/shutdown; aligns with message passing. \important{Limit}: branching into genuinely different per-UE tasks stretches the pattern.
\end{example2}
